\documentclass[12pt, letterpaper]{article}
\usepackage{graphicx}
\usepackage{amsmath}
\usepackage{amssymb}
\usepackage{float}
\graphicspath{{images/}}
\title{Normalizing features}
\author{Evert Jan Karman}
\date{January 1900}
\begin{document}
\maketitle

Normalizing features

\section{Normalizing features}
In this course, normalizing features specifically indicates: scaling each feature column by the L2 norm of the feature column.
So for each feature $h_{j}$ calculate:
\[
Z_{j} = \sqrt{\sum_{i=0}^{N-1} h_{j}(x_{i})^{2}}
\]

And then:
\[
\forall k \quad \underbar{h}_{j}(x_{k}) = \frac{h_{j}(x_{k})}{Z_{j}}
\]

Note that this happens on the \underline{training data}.
Later, the \underline{test data} must be scaled by the same factor $Z_{j}$, that was obtained using the training data.

\end{document}
